\subsection{问题四：真实地形下的地形自适应有限曲率模型}
\subsubsection{网格坐标、双线性插值与梯度}
附件网格在东西方向含 $201$ 个坐标，在南北方向含 $251$ 个坐标，两个方向步长均为 $0.02$ 海里，即 $\Delta=0.02\times1852=\SI{37.04}{m}$。记网格节点为 $(x_j,y_i)$，对应水深为 $D_{ij}$，其中 $j=0,\ldots,200$、$i=0,\ldots,250$，待测区域为 $\Omega=[0,7408]\times[0,9260]$。

对任意点 $(x,y)$，先定位其所在网格单元 $[x_j,x_{j+1}]\times[y_i,y_{i+1}]$，定义无量纲局部坐标 $u=(x-x_j)/(x_{j+1}-x_j)$、$v=(y-y_i)/(y_{i+1}-y_i)$。双线性插值模型为
\keyequation{eq:p4-interp}{
\begin{aligned}
D(x,y)=&(1-u)(1-v)D_{ij}+u(1-v)D_{i,j+1}\\
&+(1-u)vD_{i+1,j}+uvD_{i+1,j+1}.
\end{aligned}}
式~\eqref{eq:p4-interp} 在每个网格单元内可直接求导。设 $\Delta x=x_{j+1}-x_j$、$\Delta y=y_{i+1}-y_i$，则
\keyequation{eq:p4-gradient}{
\begin{aligned}
D_x&=\frac{(1-v)(D_{i,j+1}-D_{ij})+v(D_{i+1,j+1}-D_{i+1,j})}{\Delta x},\\
D_y&=\frac{(1-u)(D_{i+1,j}-D_{ij})+u(D_{i+1,j+1}-D_{i,j+1})}{\Delta y}.
\end{aligned}}
对位于最右或最上边界的点，将其归入最后一个网格单元并取 $u=1$ 或 $v=1$；超出 $\Omega$ 的坐标直接判为输入错误，不进行外推。

\subsubsection{曲线参数化与局部覆盖公式}
第 $k$ 条测线记为分段可微曲线 $\bm r_k(s)=(x_k(s),y_k(s))$。在任意非退化位置，其单位切向与左法向为 $\bm t=\bm r'_k/\|\bm r'_k\|$、$\bm n=(-t_y,t_x)$。局部横航向深度梯度为 $g=\nabla D\cdot\bm n$。令 $T=\tan\varphi$，由声束射线与局部切平面的交点关系可直接得到左、右水平覆盖半宽
\keyequation{eq:p4-swath}{
w_L=\frac{D T}{1+gT},\qquad
w_R=\frac{D T}{1-gT},\qquad
W=w_L+w_R.}
式~\eqref{eq:p4-swath} 比只写等效坡角更适合程序计算，因为 $g$ 的符号直接决定左右两侧的非对称性。其可行条件为 $D>0$、$1+gT>0$ 和 $1-gT>0$；若任一分母非正，说明该局部声束—地形几何发生奇异，候选路线被判为不可评估。

对于评估点 $\bm p$ 和一条直线段端点 $\bm a,\bm b$，令 $\bm v=\bm b-\bm a$，正交投影参数为 $\lambda=((\bm p-\bm a)\cdot\bm v)/(\bm v\cdot\bm v)$。只有 $0\le\lambda\le1$ 时投影落在线段内部；此时投影点 $\bm q=\bm a+\lambda\bm v$，有符号横向距离为 $d_\perp=(\bm p-\bm q)\cdot\bm n$。覆盖判据为
\keyequation{eq:p4-cover-test}{-w_L(\bm q)\le d_\perp\le w_R(\bm q).}
曲线在空间评估时被确定性采样为足够密的折线段，随后逐段使用式~\eqref{eq:p4-cover-test}；显式几何长度和曲率则仍按原始直线—Bézier 片段计算，避免把评估折线误当成最终曲线。

\subsubsection{评价指标与离散精度}
名义覆盖网格取 $N_x=50,N_y=60$，使用每个矩形单元中心作为评估点。设总点数 $N=N_xN_y$，未被任何测线覆盖的点数为 $N_0$，被至少两条线覆盖的点数为 $N_{\ge2}$，则漏测比例和多重覆盖比例分别为
\keyequation{eq:p4-area-metrics}{r_u=100\frac{N_0}{N},\qquad r_m=100\frac{N_{\ge2}}{N}.}

题目要求计算“重叠率超过 $20\%$ 部分的总长度”。本文要求测线按空间相邻顺序输入，对除第一条外的每条线沿中心线以不超过 \,\SI{150}{m} 的步长划分区间，并在区间中点的局部条带内均匀取 $31$ 个横向样点。设其中被前一条测线覆盖的样点数为 $m_{ki}$，则局部重叠分数 $\rho_{ki}=m_{ki}/31$；当 $\rho_{ki}>0.20$ 时，将该中心线区间长度 $\Delta s_{ki}$ 计入过度重叠长度：
\keyequation{eq:p4-overlap-metric}{L_{20+}=\sum_{k=2}^{K}\sum_i \Delta s_{ki}\,\mathbf 1\!\left(\rho_{ki}>0.20\right).}
式~\eqref{eq:p4-overlap-metric} 明确了计数政策：每个相邻条带对只在后一条测线上计数一次，不对同一物理重叠区域双重累计。

\subsubsection{候选 A、B 的保守二分布线}
候选 A 采用南北向全高直线。沿每条线取 $151$ 个等距纵向样点，第一条线选择满足所有样点左侧覆盖边界不越过西边界的最东位置。此后对每个候选位置 $q$，计算其与前一条线在全部纵向样点上的重叠率，取最小值 $\eta_{\min}(q)$；下一条线位置定义为满足 $\eta_{\min}(q)\ge0.10$ 的最大 $q$。该单调边界使用 $60$ 次二分搜索，初始区间为 $(q_{k-1}+10^{-6},7408)$ m。若前一条线在所有样点上的右边界已到达东边界，则停止。候选 B 使用完全相同的算法，只把方向旋转为东西向并交换横纵范围。两者的差异只来自地形方向性，而不是算法参数差异。

\begin{solvercontract}[保守直线候选的数值求解说明]
\item[输入] 双线性地形对象、开角 $120^\circ$、目标重叠率 $0.10$、区域边界及 $151$ 个沿程采样位置。
\item[搜索区间] A 的横向搜索区间为 $[0,7408]$ m，B 为 $[0,9260]$ m；每次二分从前一条线右侧 $10^{-6}$ m 开始。
\item[精度] 每条线位置执行 $60$ 次二分，理论区间宽度缩小为初始宽度的 $2^{-60}$；输出仍以双精度保存。
\item[停止条件] 前一条线在全部 $151$ 个沿程样点上的最小右覆盖边界达到区域末端时停止；最大允许 $500$ 条线。
\item[异常处理] 若紧邻前一条线的位置都不能达到 $10\%$ 重叠，报告“无可行下一条线”；若新位置增量小于 $10^{-4}$ m，报告递推停滞；若覆盖公式奇异，当前候选失败。
\end{solvercontract}

\subsubsection{候选 C 的站位自适应生成}
为减少保守直线在深水区造成的过度重叠，在 $y_s=s\times926$ m（$s=0,\ldots,10$）的 $11$ 个站位上分别运行单站最大间距算法，得到每个站位的有序位置序列 $\{x_{s,0},\ldots,x_{s,n_s-1}\}$。令所有站位中的最大条数为 $K=\max_s n_s=61$，第 $k$ 条曲线使用归一化序号 $t_k=k/(K-1)$。在站位 $s$ 上计算 $p=t_k(n_s-1)$、$m=\lfloor p\rfloor$、$\lambda=p-m$，并取
\keyequation{eq:p4-station-interp}{x_k(y_s)=(1-\lambda)x_{s,m}+\lambda x_{s,\min(m+1,n_s-1)}.}
式~\eqref{eq:p4-station-interp} 解决了不同站位测线条数不一致的问题：各站位都按相对序号对齐，而不是强行使用相同的局部编号。按 $y_s$ 顺序连接 $11$ 个站位点，得到 $61$ 条互不交换顺序的地形自适应折线。

\subsubsection{二次 Bézier 有限曲率平滑}
原候选 C 在每个站位处存在折点。设某条折线的连续三个节点为 $\bm P_{i-1},\bm P_i,\bm P_{i+1}$，相邻段长度为 $\ell_{i-1},\ell_i$，入射和出射单位向量为 $\bm u_i^-=(\bm P_i-\bm P_{i-1})/\ell_{i-1}$、$\bm u_i^+=(\bm P_{i+1}-\bm P_i)/\ell_i$。取平滑比例 $\tau=0.05$，截取距离 $\delta_i=\tau\min(\ell_{i-1},\ell_i)$，入口点和出口点为 $\bm E_i=\bm P_i-\delta_i\bm u_i^-$、$\bm X_i=\bm P_i+\delta_i\bm u_i^+$。用二次 Bézier 曲线
\keyequation{eq:p4-bezier}{
\bm B_i(t)=(1-t)^2\bm E_i+2t(1-t)\bm P_i+t^2\bm X_i,\qquad 0\le t\le1}
替代折点附近的两小段。由 $\bm B_i'(0)=2\delta_i\bm u_i^-$、$\bm B_i'(1)=2\delta_i\bm u_i^+$ 可知，曲线在两端分别与入射和出射直线同向，因而连接处切向连续。

曲率按
\keyequation{eq:p4-curvature}{
\kappa(t)=\frac{|x'(t)y''(t)-y'(t)x''(t)|}{\left(x'(t)^2+y'(t)^2\right)^{3/2}},\qquad R(t)=\frac{1}{\kappa(t)}}
计算。式~\eqref{eq:p4-curvature} 将曲线的一、二阶导数转化为可直接比较的曲率半径。每个 Bézier 段以 $401$ 个参数点检查最小曲率半径，弧长使用 $64$ 子区间复合 Simpson 积分；用于覆盖评估时，每个曲线段再均匀细分为 $4$ 段。若导数为零、相邻片段不连通或连接方向差超过 $10^{-6}$ 度，则判为几何失败。

\begin{modelsummary}[问题四模型归纳]
问题四模型由四个互相独立但可衔接的模块组成：地形模块用式~\eqref{eq:p4-interp}--\eqref{eq:p4-gradient} 给出任意位置水深与梯度；条带模块用式~\eqref{eq:p4-swath} 和~\eqref{eq:p4-cover-test} 判断空间覆盖；路线模块用保守二分、站位归一化插值和式~\eqref{eq:p4-bezier} 生成显式有限曲率曲线；评估模块用式~\eqref{eq:p4-area-metrics} 和~\eqref{eq:p4-overlap-metric} 输出题目规定指标。生成器不读取评估结果反向修改指标定义，评估器也不负责选择候选。
\end{modelsummary}

\subsubsection{完整求解算法与特殊情况处理}
问题四的算法顺序如下：首先读取并检查网格坐标严格递增、水深均为正且无缺失；其次构造双线性地形对象；第三步按统一参数生成 A、B 和 C；第四步对 C 使用 $\tau\in\{0.05,0.10,0.20,0.30,0.40\}$ 生成五个显式平滑变体；第五步检查边界、片段连通性、硬拐角、曲率、相邻线交叉和最小间距；第六步使用冻结的名义网格和重叠采样参数独立评估；第七步按漏测率、过度重叠长度、总长度的字典序比较；第八步在 $30\times30$、$50\times50$、$70\times70$ 网格上进行敏感性复算并记录不利结果。

\begin{solvercontract}[问题四算法求解说明]
\item[输入] $251\times201$ 正水深网格，区域 $[0,7408]\times[0,9260]$ m，开角 $120^\circ$，目标局部重叠率 $0.10$。
\item[生成精度] 直线位置二分 $60$ 次；A、B 的沿程最不利值各取 $151$ 个样点；C 使用 $11$ 个站位；平滑比例候选为 $0.05$ 至 $0.40$。
\item[评估精度] 名义覆盖网格为 $50\times60$；中心线步长不超过 $150$ m；每个局部条带取 $31$ 个横向样点；敏感性网格为 $30,50,70$。
\item[可行性条件] 水深和条带分母必须为正；路线不得越界；相邻线不得交叉；显式几何必须连续且无硬拐角；曲线导数不得为零。
\item[特殊情况] 网格边界采用最后单元插值；不同站位条数通过归一化序号插值；没有船舶最小转弯半径时只报告 $R_{\min}$；高分辨率网格出现漏测时保留原结果并降低结论强度，不事后修改阈值。
\item[输出] 每条线的 $11$ 站位坐标、显式几何长度和最小曲率半径，以及方案级条数、总长度、漏测比例、多重覆盖比例、$L_{20+}$、越界长度和敏感性结果。
\end{solvercontract}

\subsubsection{候选结果比较}
表~\ref{tab:p4-compare} 汇总四类主要候选。A 由 $67$ 条全高直线组成，因此总长度严格为 $67\times9260=\SI{620420}{m}$；B 由 $95$ 条全宽直线组成，总长度为 $95\times7408=\SI{703760}{m}$。B 的长度和过度重叠均大于 A，说明在该地形上东西向全局路线不具优势。原 C 仅有 $61$ 条，长度显著下降，但 $61$ 条线的 $9$ 个内部站位共形成 $61\times9=549$ 个潜在折点，其中一个折点方向变化低于数值阈值，最终检测到 $548$ 个硬拐角，因此不能直接作为作业路线。C-05 经显式平滑后硬拐角为零，且长度和 $L_{20+}$ 都略低于原 C。

\begin{table}[H]
\centering
\caption{问题四候选方案比较}
\label{tab:p4-compare}
\begin{tabularx}{\textwidth}{C{1.5cm}C{1.3cm}C{2.45cm}C{2.45cm}C{1.8cm}L}
\toprule
方案 & 条数 & 总长度/m & $L_{20+}$/m & 名义漏测/\% & 可行性与角色 \\
\midrule
A & 67 & 620420.000 & 369055.806 & 0 & 全部声明的敏感性网格均为零漏测，保守回退 \\
B & 95 & 703760.000 & 452776.960 & 0 & 可行，但在长度和过度重叠上被 A 支配 \\
C & 61 & 576755.702 & 339673.098 & 0 & 含 $548$ 个硬拐角，不作为最终路线 \\
C-05 & 61 & 576726.636 & 339556.407 & 0 & 切向连续、无交叉、无边界违规，最终推荐 \\
\bottomrule
\end{tabularx}
\end{table}

图~\ref{fig:p4-compare} 显示，C-05 相对 A 的总长度减少 $7.043\%$，$L_{20+}$ 减少 $7.993\%$。这种改善来自测线密度随局部地形调整，而不是通过放宽名义覆盖要求获得。多重覆盖面积比例仍约为 $47.57\%$，说明复杂地形下局部过度覆盖尚未完全消除，后续改进应重点处理浅水突变和曲线边界附近，而不是继续压缩整体条数。

\begin{figure}[H]
\centering
\includegraphics[width=0.72\textwidth]{problem4_comparison.pdf}
\caption{问题四候选方案总长度与过度重叠长度比较}
\label{fig:p4-compare}
\end{figure}

\subsubsection{最终方案与题目指标的直接回答}
最终推荐 C-05，具体路线见图~\ref{fig:p4-routes}。该方案包含 $61$ 条由直线段与二次 Bézier 段组成的显式曲线，全部位于区域内，相邻线顺序不交换。其几何总长度为 \,\SI{576726.636}{m}，独立空间评估器对离散折线复算得到 \,\SI{576725.729}{m}，两者相差 \,\SI{0.906}{m}，相对差约 $1.57\times10^{-6}$。名义 $50\times60$ 网格中没有未覆盖单元中心，因此题目要求的漏测海区占比在该评估口径下为 $0\%$；按式~\eqref{eq:p4-overlap-metric} 得到重叠率超过 $20\%$ 部分的总长度为 \,\SI{339556.407}{m}。这三个数值分别直接回答题目四要求的总长度、漏测比例和过度重叠长度。

\begin{figure}[H]
\centering
\includegraphics[width=0.68\textwidth]{problem4_routes.pdf}
\caption{问题四最终推荐的 $61$ 条地形自适应有限曲率测线}
\label{fig:p4-routes}
\end{figure}

图~\ref{fig:p4-curvature} 给出各测线最小曲率半径。方案整体最小值为 \,\SI{225.983}{m}，硬拐角数为 $0$，相邻线交叉数为 $0$，最小相邻间距为 \,\SI{59.474}{m}。这些结果证明路线几何连续且曲率有限，但因为题目没有船型参数，不能据此推出船舶一定能够以给定航速完成转弯。

\begin{figure}[H]
\centering
\includegraphics[width=0.70\textwidth]{problem4_curvature.pdf}
\caption{最终方案各测线的最小曲率半径}
\label{fig:p4-curvature}
\end{figure}
